\section{Results}
\label{sec:battery_results}

    \subsection[Assessing e-bikes' battery levels and usage patterns]{Assessing \acp{e-b}' battery levels and usage patterns}
    \label{subsec:battery_results_analysis}

        The average battery levels per neighborhood show significant variability across neighborhoods (Figure~\ref{battery_overiew}A, Appendix~\ref{app_paper2:subsection_battery_levels_across_time}). To facilitate the analysis of all the temporal series of average battery levels derived from each neighborhood, K-means clustering was performed, resulting in four distinct clusters (Figure~\ref{battery_overiew}B, Appendix~\ref{app_paper2:subsection_choice_of_number_of_clusters}). The central and southern neighborhoods, which fall into Clusters 1 and 4, have the lowest battery levels, exhibiting steep declines in battery levels during business hours with recovery during nighttime (Figure~\ref{battery_overiew}C). This is likely due to the high levels of industrial and commercial activity, as well as the dense populations in these areas. Additionally, it is worth noting that some neighborhoods in Cluster 4 have slower charging speeds at night, likely due to a stronger nightlife presence. In contrast, northern neighborhoods in Clusters 2 and 3 are mainly residential with less commercial activity, resulting in higher and more stable battery levels throughout the day.

        \begin{figure}[!htbp]
          \centering
          \includegraphics[width=0.9\textwidth]{0_plots/paper2/main/Panel1.png}
          \caption[Average neighborhood battery levels and temporal clustering]{
            \textbf{Average neighborhood battery levels and temporal clustering.}
            \textbf{(A)} Average battery levels across neighborhoods. The average for each neighborhood is computed using the battery values of all the bikes in that neighborhood. 
            \textbf{(B)} Clustering of the neighborhoods according to the temporal pattern of battery levels using K-means into four distinct clusters. 
            \textbf{(C)} Temporal series of average battery levels for each cluster. The neighborhoods in grey do not have bike-sharing stations.}
          \label{battery_overiew}
        \end{figure}

        One factor directly related to average battery levels is the average time between trips for each bike, which can also be referred to as inter-event time. With longer inter-event times, \acp{e-b} have more opportunity to recharge, resulting in higher battery levels. In Figure~\ref{correlations}A, central and southern neighborhoods (Clusters 1 and 4) have not only lower average battery levels but also shorter inter-event times, indicating a higher turnover of \acp{e-b} due to frequent trips. This high turnover aligns with the higher density of users and increased mobility demand in commercial and densely populated areas. However, Zona Franca, located in the southernmost area of the city, presents a distinct case due to its industrial nature and significantly low population density (1 inhabitant per hectare \cite{bcn_population_density}), which results in low \ac{e-b} usage. Conversely, northern neighborhoods (Clusters 2 and 3) exhibit longer inter-event times, reflecting less frequent \ac{e-b} usage and consequently higher and more stable battery levels despite some of them being in high altitudes. Additionally, a Pearson correlation analysis was conducted between average battery levels and inter-event time at the station level (Figure~\ref{correlations}B), revealing a correlation value of 0.39. This moderate positive correlation suggests that stations with longer inter-event times tend to have higher average battery levels, further supporting the neighborhood-level results.

        \begin{figure}[!htbp]
            \centering
            \includegraphics[width=0.9\textwidth]{0_plots/paper2/main/Panel2.png}
            \caption[Inter-event times, distance index, and their relationships with battery levels]{
              \textbf{Inter-event times, distance index, and their relationships with battery levels.} 
              \textbf{(A)} Average inter-event time between trips across neighborhoods. The average for each neighborhood is computed using the inter-event time values of all bikes, with outliers removed using the IQR method at the station level. 
              \textbf{(B)} Relation between average battery levels and inter-event time at the station level, with a Pearson correlation of 0.39 (p-value<0.001). 
              \textbf{(C)} Distance index per neighborhood, calculated by first computing the index at the station level to accurately reflect trip distances between stations, and then averaging it at the neighborhood level. 
              \textbf{(D)} Relation between average battery levels and the distance index at the station level, with a Pearson correlation of 0.47 (p-value<0.001). The neighborhoods in gray do not have bike-sharing stations.}
            \label{correlations}
        \end{figure}

        Apart from the inter-event time, two additional factors can affect the average battery levels at a station: the origins of the trips that the station receives and the number of trips coming from each origin. To evaluate the impact of these factors on battery depletion, a distance index (\(DI_j\)) was developed for each station \( j \). This index measures the influence of the distances traveled by bikes on their battery levels by accounting for both the distance of each trip and the frequency of trips from each origin station. Specifically, for trips originating at station \( i \) and ending at station \( j \), the distance  (\( d_{i,j} \)) is weighted by the inverse of the total trips originating from station \( i \) (\( n_i \)). This weighting ensures that trips from less frequent origins have a greater influence on the distance index. The weighted distances for all origin stations \( i \) are then summed for the destination station \( j \). To make the index independent of the station’s overall demand, the sum of these weighted distances is normalized by the total number of trips received by station \( j \) (\( \sum_{i} T_{ij} \)). The resulting distance index provides a measure of how trip distances, adjusted for trip frequency, contribute to battery depletion at a station. In intuitive terms, higher \(DI_j\) means the station receives fewer trips overall and a greater share from distant origins, so bikes there spend more time docked recharging and less time in use. Because trips drain the battery while docking recharges it, this recharge effect dominates and higher \(DI_j\) is therefore associated with higher average battery levels. The formula of the index is as follows:

        \begin{equation}
          DI_j = \frac{\sum_{i} \frac{d_{i,j}}{n_i}}{\sum_{i} T_{ij}}
        \end{equation}

        The distance index was computed for all stations, and averaged over each neighborhood (Figure~\ref{correlations}C). As shown, neighborhoods in the northern areas tend to have a higher distance index, indicating that trips to these neighborhoods often originate from farther stations and that there are fewer incoming trips overall. This combination results in higher average battery levels. Conversely, neighborhoods in the central and southern parts of the city exhibit a lower distance index, suggesting that trips to these areas typically come from nearby stations and that there are more incoming trips, leading to lower battery levels. Zona Franca, as previously mentioned, deviates from this trend due to its unique industrial nature and low population density. The Pearson correlation between the distance index and average battery levels at the station level revealed a value of 0.47 (Figure~\ref{correlations}D), indicating a moderate positive correlation. This suggests that stations farther from others and receiving a lower number of trips tend to have higher average battery levels. Additionally, no significant collinearities were found between the inter-event time and the distance index (Appendix~\ref{app_paper2:subsection_collinearity}).


    \subsection[Modeling e-bikes' battery levels]{Modeling \acp{e-b}' battery levels}

        Using the battery model described in Section~\ref{sec:battery_modeling} (linear increase when docked, decrease during trips), the optimal parameters for the power consumption equation and charging speed were determined through hyperparameter optimization (see Section~\ref{sec:battery_optimization}). The fit between the observed and expected battery levels, using these optimal parameters, is shown in Figure \ref{fig:battery_fit}. In Figure \ref{fig:battery_fit}A, the predicted battery percentage is displayed, while in Figure \ref{fig:battery_fit}B, the predicted battery variation is presented. Despite the dispersion, a significant correlation is observed, with values closely aligned to the diagonal. Factors such as user heterogeneity, missing route information, or bikes not charging while docked could contribute to the dispersion. The parameters were also validated using the 8-hour dataset, with similar correlation values (Appendix~\ref{app_paper2:subsection_linear_battery_function}). In addition, a non-linear charging function was tested using a sigmoid-like function to account for faster charging times at lower battery levels. The fitting results are very similar to the linear function, with the optimal shape inferred being almost linear (Appendix~\ref{app_paper2:subsection_nonlinear_charging_function}).

        \begin{figure}[!htbp]
          \begin{center}
          \includegraphics[width=0.9\textwidth]{0_plots/paper2/main/battery_prediction.png}
          \end{center}
          \caption[Battery model validation using the 30-minute battery dataset]{
            \textbf{Battery model validation using the 30-minute battery dataset.} 
            \textbf{(A)} Observed battery level as a function of the predicted values. 
            \textbf{(B)} Battery variation observed as a function of the predicted values. Each point corresponds to one battery-level observation. The Root Mean Squared Error (RMSE) is $0.0178$ and the Pearson correlation coefficient is significant in both cases (p-val<0.001). The linear fit is shown in gray.
          } 
          \label{fig:battery_fit}
        \end{figure}

        This battery model can have two key applications. The first enables the assessment of battery consumption and destination probabilities for individual trips. The second integrates the Markov chain-based approach with the battery model to predict both \ac{e-b} locations and battery levels across the entire fleet.


    \subsection{Single trip application}

        Evaluating short-term scenarios for a single bike can be achieved by coupling the battery model with the distribution of possible destinations $P_{ij}^{w,h}$. In Figure \ref{fig:single_trip}, given \textit{St. Ciutat de Granada, 168} as the origin station, the probability of arriving at each destination is represented by the size of the dots, while the battery consumption for each bike trip is indicated by the color. The probability of destinations changes significantly between 8:00 and 18:00 despite the same origin station. At 18:00, trips tend to be shorter and head towards nearby stations, while morning trips are longer and converge towards a few work-related hot spots. Additionally, battery consumption varies notably based on the elevation changes along the route, with trips going downhill consuming less battery over the same distance. 

        \begin{figure}[!htbp]
          \begin{center}
          \includegraphics[width=0.9\textwidth]{0_plots/paper2/main/single_trip_daytype=0_station=42.pdf}
          \end{center}
          \caption[Single trip location and battery prediction]{
            \textbf{Single trip location and battery prediction.}
            Probability of finding a bike at a station and battery consumption of the trip when departing from the station in \textit{St. Ciutat de Granada, 168} on a weekday at \textbf{(A)} 08:00 and \textbf{(B)} 18:00. The origin station is highlighted with a triangle. The size of the dots corresponds to the probability of having that station as the destination and color corresponds to the battery consumed by the bike. The grey points correspond to destination stations with no trips.
          }
          \label{fig:single_trip}
        \end{figure}

        Furthermore, an impact index can be computed to assess the influence of each origin station on battery levels. This index depends on the probability of each destination given an origin station, the battery consumption of each possible trip, and the average hourly trips departing from station $i$ to other stations. The impact index can be written as:

        \begin{equation}
        \gamma^{w,h}_i=\langle T^{w,h}_i \rangle P^{w,h}_{ij} \Delta B_{ij}
        \end{equation} 

        where $\langle T^{w,h}_i \rangle$ represents the average number of trips departing from station $i$, $P^{w,h}_{ij}$ is the probability that a trip generated in station $i$ ends in station $j$, and $\Delta B_{ij}$ is the battery consumed by a trip from $i$ to $j$. An example of the impact index $\gamma$ at 8:00 and 18:00 on a regular weekday is shown in Figure \ref{fig:impact}. Due to the variation in trip patterns captured by $P^{w,h}_{ij}$, the stations with the highest impact change throughout the day.

        \begin{figure}[!htbp]
          \begin{center}
          \includegraphics[width=0.9\textwidth]{0_plots/paper2/main/impact_index_0.pdf}
          \end{center}
          \caption[Station impact index at different hours of the day]{
            \textbf{Station impact index at different hours of the day.} Impact index of the stations at \textbf{(A)} 08:00 and \textbf{(B)} 18:00. The station with higher impact at each hour is highlighted with a larger dot size. In the morning, the \textit{St. Provença, 125} station has the highest impact, while in the afternoon, \textit{St. Ciutat de Granada, 168} has the greatest impact.
            } 
          \label{fig:impact}
        \end{figure}

    \subsection{Coupling the Markov chain and the battery models}

        The adjusted battery consumption function can be integrated into the Markov chain approach to provide joint predictions of both location and battery levels over multiple trips, either at the single-bike or \ac{bss} level. Inter-event times are used to calculate the charging percentage, while the distance and duration of trips are used to calculate energy consumption. Thus, the inputs required for the coupled model are the location and battery level of each bike along with the day of the week and the hour, which together define the starting point of the simulation.

        To validate this approach, systematic simulations were performed using initial conditions from the 2023 mobility data and 30-minute battery data (Figure \ref{fig:markov_correlations}). Due to the inherent noise in micro-mobility service usage (i.e., the unpredictability of which bike at a station will be selected by the next user), the validation was conducted at the station level. For each validation snapshot, the initial conditions were determined by the number of available bikes and their average battery level at each station in the snapshot. To address potential gaps in battery data, the battery level of each bike was set to the station average. However, while comparisons between actual and predicted values were made at the station level, the mobility and battery dynamics of each bikes were simulated individually. Additionally, to prevent data leakage, the probabilities of the Markov model were based on data from the year 2022.

        \begin{figure}[!htbp]
          \begin{center}
          \includegraphics[width=0.9\textwidth]{0_plots/paper2/main/markov_correlations.pdf}
          \end{center}
          \caption[Correlations for the Markov chain model across prediction horizons.]{
            \textbf{Correlations for the Markov chain model across prediction horizons.} 
            \textbf{(A)} Pearson correlation coefficient between actual and predicted bike availability per station as a function of $\Delta t$. 
            \textbf{(B)} Pearson correlation coefficient between actual and predicted average battery level per station as a function of $\Delta t$.
          }
          \label{fig:markov_correlations}
        \end{figure}

        Once both the mobility and battery consumption models were validated, simulations across various initial scenarios and time horizons were conducted to demonstrate the potential applications of our framework (Appendix~\ref{app_paper2:section_initial_conditions_for_battery_predictions}). For these predictions, the initial conditions were derived from the average number of available bikes and their battery levels across stations in snapshots corresponding to the same type of day and hour. First, it is possible to compute the probability of finding a bike at a given destination and its battery level as a function of its origin, following one or more trips within a time period $\Delta t$. The final battery levels are determined by averaging over all the Markov iterations. As an example, Figure \ref{fig:markov_trips} shows the probability and battery levels of bikes departing from the station at Glòries (\textit{St. Ciutat de Granada, 168}) at 08:00 after \textbf{(A)} 30 and \textbf{(B)} 120 minutes. The role of distance from the origin is evident in the short-term predictions, with lower battery levels at greater distances and inclinations. However, this effect diminishes over longer time windows, as multiple trips can occur before reaching a station. For longer periods, a greater variety of final destinations is also possible.


        \begin{figure}[!htbp]
          \begin{center}
          \includegraphics[width=0.9\textwidth]{0_plots/paper2/main/markov_trips_daytype=0-hour=08-station=42.pdf}
          \end{center}
          \caption[Predicted bike locations and battery levels for trips starting at 8:00 from Glòries station]{
            \textbf{Predicted bike locations and battery levels for trips starting at 8:00 from Glòries station.} Probability of finding a bike and average battery levels after \textbf{(A)} 30 minutes and \textbf{(B)} 120 minutes of simulation for trips departing from the Glòries station (St.\ Ciutat de Granada, 168) at 08:00 on an average weekday. The origin station is highlighted with a triangle.
          }
          \label{fig:markov_trips}
        \end{figure}

        As a second application, the Markov chain approach allows to transition from bike-level to system-level predictions. To calculate the number of available bikes per station across the entire fleet, the probabilities $f^{w,h}_{i,j}(\Delta t)$ for each destination station $j$ are summed, which may result in non-integer values. Battery levels, on the other hand, are calculated by multiplying each bike's probability by its corresponding battery level and then dividing by the number of available bikes. Figure~\ref{fig:markov_battery} displays the number of available bikes and the average battery level after two hours of simulation for two different starting times, 08:00 and 18:00. Stations near the city center exhibit lower battery levels, consistent with the data analysis results. Meanwhile, the northern part of the city, which experiences lower usage rates, shows higher battery levels. Additionally, lower battery levels are observed in the afternoon due to increased battery consumption throughout the day.

        \begin{figure}[!htbp]
          \begin{center}
          \includegraphics[width=0.9\textwidth]{0_plots/paper2/main/markov_battery_daytype=0-dt=120.pdf}
          \end{center}
          \caption[System-level predictions of bike availability and battery levels after 120 minutes of simulation]{
            \textbf{System-level predictions of bike availability and battery levels after 120 minutes of simulation.} 
            Predictions combine the Markov mobility model with the battery equations for simulations starting at \textbf{(A)} 08:00 and \textbf{(B)} 18:00 on an average weekday.
          }
          \label{fig:markov_battery}
        \end{figure}
